\section{Fully charged recursive schedules}\label{sec:schedules}

All size parameters in this section refer to the original instance. In
particular, we do not restart a low-degree schedule at a child whose
tensors were created at an earlier depth. Write
\[
 \lambda=\ceil{\log_2s},\quad
 \Lambda=\lambda-\ceil{\sqrt\lambda},\quad
 b=\ceil{6\log_2\lambda},\quad H=\ceil{s/\lambda}.
\]
Directly contract instances with at most $H$ wires. For all larger calls
at global depth $d$, choose
\begin{equation}\label{eq:integer-schedule}
 u_d=\floor{\eta_d\Lambda},\qquad
 a_d=\sum_{j=0}^d u_j+(d+1)b.
\end{equation}
The finitely many sizes with nonpositive tapers are also treated directly.
All these choices define one uniform algorithm after fixing the backend
accuracy. Rounding the logarithms can be implemented with integer bit
lengths and integer power comparisons.

Both schedules have four initial stages followed by a four-group tail.
Set $q=63/100$. If $A_*$ denotes the total taper budget and
$A_3=\sum_{j=0}^3\eta_j$, then
\begin{equation}\label{eq:tail-taper}
 r_d=4,\qquad \eta_d=(A_*-A_3)(1-q)q^{d-4}\quad(d\ge4).
\end{equation}

\begin{center}\small
\begin{tabular}{@{}crrrr@{}}
\toprule
 & \multicolumn{2}{c}{Generic root degree $s^{9/10}$}
 & \multicolumn{2}{c}{Low root degree $s^{1/12}$}\\
Stage $d$ & $r_d$ & $\eta_d$ & $r_d$ & $\eta_d$\\\midrule
0 & 10 & $301/20000$ & 15 & $2149/20000$\\
1 & 13 & $287/25000$ & 10 & $2337/100000$\\
2 & 10 & $357/100000$ & 7 & $407/20000$\\
3 & 7 & $193/50000$ & 6 & $109/10000$\\
4 & 4 & $14837/2500000$ & 4 & $140341/10000000$\\\midrule
Total $A_*$ & \multicolumn{2}{c}{$1/20$}
 & \multicolumn{2}{c}{$1/5$}\\
Tail budget $A_*-A_3$ & \multicolumn{2}{c}{$401/25000$}
 & \multicolumn{2}{c}{$3793/100000$}\\
\bottomrule
\end{tabular}
\end{center}

\subsection{Packing gates and a uniform depth argument}

Put $A_d=\sum_{j\le d}\eta_j$ and $A_{-1}=0$. If the original degree
exponent is $\delta$, the root has both side degrees at most
$s^\delta$ up to rounding. Below the root the fresh-side degree is at
most $2^{a_{d-1}+1}$; a safe retained-side degree bound is the maximum
of the original degree and all earlier packet degrees. For any fixed
prefix depth, the leading exponents in the numerators of
\eqref{eq:gates} are therefore bounded by
\[
\begin{array}{c|ccc}
 & QR & Q^2C & L_pQ^2\\\hline
 d=0 & \delta+A_0\log_2r_0
       &\delta+2A_0\log_2r_0&A_0(1+2\log_2r_0)\\
 d>0 & \max(\delta,A_{d-1})+A_d\log_2r_d
       &A_{d-1}+2A_d\log_2r_d&A_d(1+2\log_2r_d).
\end{array}
\]
Since an active call has $M>s/\lambda$, each strict bound below one
establishes its gate, including any fixed polylogarithmic factors.

An independent rational audit certifies the following upper bounds. The
displayed terminating decimals are exact rational upper bounds; entries
with ellipses are shown rounded for readability.

\begin{center}\small
\begin{tabular}{@{}ccrrr@{}}
\toprule Schedule & $d$ & $QR$ & $Q^2C$ & $L_pQ^2$\\\midrule
Generic &0& .9499961&.9999922&.1150422\\
 &1&.998174265&.21139853&.22287853\\
 &2&.9999922&.2265144&.2300844\\
 &3&.995339304&.220778608&.224638608\\\midrule
Low degree &0&.5031404833\ldots&.9229476333\ldots&.9470643\\
 &1&.54203404&.97661808&.99998808\\
 &2&.555214658&.979609316&.999959316\\
 &3&.57012095&.9890719&.9999719\\
\bottomrule
\end{tabular}
\end{center}

For example, the audit verifies $r^v<2^u$ exactly to prove the rational
logarithm bounds $\log_2r<u/v$. The smallest certified prefix gaps are
$39/5000000$ and $149/12500000$, respectively. They are positive fixed
margins, though their small sizes permit extremely large asymptotic
thresholds. No practical crossover claim follows.

There are boundary cases in the tail, and ignoring them would be an
error. In the generic case the largest leading exponent is
$9/10+2A_*=1$ in $QR$. The use of $\Lambda$ instead of $\lambda$
gives the sharper bound
\[
 \log_2(QR)\le\lambda-\tfrac1{10}\sqrt\lambda
                        +O((\log\lambda)^2).
\]
The other two numerator exponents are at most $1/4+o(1)$.
In the low-degree case, $QR$ has exponent at most $3/5+o(1)$, while
both remaining gates have boundary exponent $5A_*=1$ and satisfy
\[
 \max\{\log_2(Q^2C),\log_2(L_pQ^2)\}
 \le\lambda-\sqrt\lambda+O((\log\lambda)^2).
\]
The negative square-root margins dominate division by $M$, the required
logarithmic factors, and accumulated integer guards.

To justify the depth used in these estimates without circular reasoning,
first consider hypothetical active depths up to
\[
 D_\lambda=\ceil{\log_{20/11}\lambda}+O(1).
\]
On that whole range, $d=O(\log\lambda)$, so all guards total
$O((\log\lambda)^2)$. The smallest tail taper is at least a constant
times
\[
 \lambda^{1+\ln(63/100)/\ln(20/11)},
 \qquad 1+\frac{\ln(63/100)}{\ln(20/11)}>0.
\]
Positivity follows exactly from $63/100>11/20$. Thus every active taper
tends to infinity uniformly. The already established gates now imply
tail child ratio $1/2+o(1)\le11/20$ for all sufficiently large $s$.
Such an active path cannot exceed $D_\lambda$, because its size would
then fall below $s/\lambda$. This closes the induction. Use relative
packing slack $\lambda^{-3}$; over $O(\log\lambda)$ levels its
accumulation is $1+o(1)$. Floors are uniformly negligible.

\subsection{Precision and the costs outside the main exponent}

Set $a_{-1}=0$. The inherited precision satisfies
$p=O(s\lambda^2 2^{a_{d-1}})$. For a retained valuation $i<u_d$,
\[
 D=2^{a_d-i}\ge2^{a_{d-1}+b+1},\qquad
 \frac{p}{MD/u_d}=O(\lambda^4 2^{-b})=o(1).
\]
The exact Hermite slope is $(1+o(1))M/u_d$, and its remaining allowance
has relative size $O(\lambda^2/D)=o(1)$. This proves the uniform sharp
working-precision estimate $(6+o(1))MD/u_d$ used in the entropy bound.
The optimal mode tilt has scale $\sqrt{u_d}/D\to0$ uniformly, making
the generating-function remainder negligible as well.

Throughout the recursion,
$a_d\le A_*\lambda+O((\log\lambda)^2)$, so packet lengths are at
most $s^{A_*+o(1)}$. Mode, precision, interpolation, and packet
descriptions remain polynomial in $s$. The source extension degree
$4\ceil{\log_2s}$ supplies enough residue-distinct points, because
$M(2L_p-1)+1<s^4$ eventually. Entry queries retain their creation
precision. Fresh packet evaluation costs
$2^{O(L_p\log s)}=2^{s^{A_*+o(1)}}$, original entry work costs
$2^{O(s^{9/10})}$, and direct leaves cost $2^{O(s/\lambda)}$.
Each is $2^{o(s/\sqrt\lambda)}$ under both schedules. Products of
polynomial factors over the depth add only $O(\lambda\log\lambda)$
to the exponent. Description-dependent work can be charged once per
call; the branching bounds do not depend on $\LL$. The dependence on
the original description therefore remains one fixed polynomial factor.

\subsection{Exact global coefficients and the terminal route}

Let $\mu_0=1$, $\mu_{d+1}=2\mu_d/r_d$, and
\[
 k_d=\kappa\mu_d/\sqrt{\eta_d},\qquad
 P_d=\sum_{j<d}\alpha_{r_j}k_j.
\]
Unrolling the paid recurrence \eqref{eq:paid}, the local route at depth
$d$ has coefficient $P_d+\beta_{r_d}k_d$; its own table space has
coefficient $\beta_{r_d}k_d$. The wire masses through depth four are
\[
 \begin{array}{c|ccccc}
 &0&1&2&3&4\\\hline
 \text{generic}&1&1/5&2/65&2/325&4/2275\\
 \text{low degree}&1&2/15&2/75&4/525&4/1575.
 \end{array}
\]
In the tail, $k_{4+j}=k_4\gamma^j$ with
$\gamma=1/(2\sqrt q)=5/\sqrt{63}<1$. Since
$\beta_4<\alpha_4/(1-\gamma)$, the largest possible tail route is
its limiting terminal coefficient
\begin{equation}\label{eq:terminal}
 P_4+\frac{\alpha_4 k_4}{1-\gamma},\qquad
 \alpha_4=\sqrt{3/8}.
\end{equation}
It pays for every recursive table entry down to the direct leaves. For
finite tail depth this is an upper bound. The unspent tail mass tends to
zero, so it also specifies the leading limiting coefficient.

\begin{center}\small
\begin{tabular}{@{}lrr@{}}
\toprule Charged route & Generic & Low degree\\\midrule
Local stage 0 &28.632766276536&10.336654568635\\
Local stage 1 &28.585713591899&9.961674795027\\
Local stage 2 &28.484720584611&10.108413259516\\
Local stage 3 &28.520568607891&10.174645147995\\
First tail stage &28.493870083097&10.206500056211\\
Terminal route &28.632507095859&10.336723985547\\
\bottomrule
\end{tabular}
\end{center}

The generic maximum is the root. The low-degree maximum is the terminal
route, slightly larger than its root. Accordingly the exact constants are
\begin{align*}
 C_{\rm gen}&=\kappa\beta_{10}\sqrt{20000/301},\\
 C_{\rm low}&=\kappa\left(
  \sum_{d=0}^3\frac{\alpha_{r_d}\mu_d}{\sqrt{\eta_d}}
  +\frac{\alpha_4\mu_4}{\sqrt{\eta_4}(1-5/\sqrt{63})}\right),\\
 C_{\rm low,space}&=\kappa\beta_{15}\sqrt{20000/2149}.
\end{align*}
The middle expression uses the low-degree row of the parameter table.
Rational interval arithmetic proves $C_{\rm gen}<28.633$ and
$C_{\rm low}<10.337$, as well as the indicated branch dominance.
The intervals use an alternating arctangent series for $\pi$, the
positive atanh series for $\ln2$, and rational square-root enclosures.
No floating-point value decides a claimed inequality.

Ancestor tables are kept while a child is evaluated. Their space adds,
rather than multiplies, so the exponent coefficient is the largest
local table coefficient over the recursion. In both schedules the root
is largest; the decreasing geometric tail and the finitely many prefix
comparisons certify this. Polynomial local workspace and recursion
descriptions do not alter the coefficient. This completes the proof of
Theorem~\ref{thm:contractors} relative to the explicitly imported interface.
